\documentclass[11pt]{article}
\input{style-sujet}
\usepackage{array}

\begin{document}

\entetesujet{Sujet bac 2013 -- Série D}

% ====================== EXERCICE 1 ======================
\exo{1}{4 points}

Les caractères $X$ et $Y$ sont distribués suivant le tableau à double entrée ci-après :

\begin{center}
\renewcommand{\arraystretch}{1.3}
\begin{tabular}{|c|c|c|c|}
\hline
$X\,\backslash\,Y$ & $-1$ & $0$ & $2$ \\
\hline
$-2$ & 4 & 0 & 2 \\
\hline
$-1$ & 3 & 5 & 0 \\
\hline
$0$ & 2 & 1 & 2 \\
\hline
\end{tabular}
\end{center}

\begin{enumerate}
  \item Dresser la loi marginale de $X$ et celle de $Y$.
  \item Trouver les coordonnées du point moyen $G(\overline{X},\overline{Y})$.
  \item Déterminer l'équation de la droite de régression de $Y$ en $X$.
  \item Calculer le coefficient de corrélation linéaire de la série statistique.
\end{enumerate}

% ====================== EXERCICE 2 ======================
\exo{2}{4 points}

Le plan $(P)$ est rapporté à un repère orthonormé $(O,\vec{u},\vec{v})$.

\begin{enumerate}
  \item Résoudre dans l'ensemble des nombres complexes, l'équation $(E) : Z^3 + 8 = 0$. On donnera les solutions de $(E)$ sous la forme algébrique.
\end{enumerate}

Soit $A$, $B$, et $C$ les points d'affixes des complexes : $Z_A = 1 + i\sqrt{3}$ ; $Z_B = -2$ ; $Z_C = 1 - i\sqrt{3}$.

\begin{enumerate}[start=2]
  \item \begin{enumerate}[label=\alph*.]
    \item Calculer le module et un argument de $U$ tel que : $U = \dfrac{Z_C - Z_A}{Z_B - Z_A}$.
    \item En déduire la nature du triangle $ABC$.
  \end{enumerate}
\end{enumerate}

Soit $S$, la rotation définie dans $(P)$ telle que : $S(A) = C$ et $S(C) = B$.

\begin{enumerate}[start=3]
  \item \begin{enumerate}[label=\alph*.]
    \item Déterminer l'expression complexe de $S$.
    \item Déterminer les éléments caractéristiques de $S$.
  \end{enumerate}
\end{enumerate}

% ====================== PROBLÈME ======================
\prob{12 points}

\partie{A}
Soit $g$ la fonction numérique de la variable réelle $x$ définie sur $]0\,;+\infty[$ par :
\[ g(x) = 1 - x^2 - \ln x \]

\begin{enumerate}
  \item Étudier les variations de $g$, puis dresser son tableau de variation.
  \item Calculer $g(1)$, puis en déduire le signe de $g$ sur $]0\,;+\infty[$.
\end{enumerate}

\partie{B}
On considère la fonction numérique $f$ de la variable réelle $x$ définie par :
\[ f(x) = \begin{cases} 1 + x - e^{1-x} & \text{si } x \le 1 \\[1ex] \dfrac{2x - x^2 + \ln x}{x} & \text{si } x > 1 \end{cases} \]
On désigne par $(C)$, la courbe représentative de $f$ dans le repère orthonormé $(O,\vi,\vj)$ du plan d'unité graphique : 2 cm.

\begin{enumerate}
  \item Déterminer l'ensemble de définition de la fonction $f$.
  \item \begin{enumerate}[label=\alph*.]
    \item Étudier la continuité et la dérivabilité de $f$ au point $x = 1$.
    \item Pour $x \in ]1\,;+\infty[$, exprimer $f'(x)$ en fonction de $g(x)$.
  \end{enumerate}
  \item Étudier les variations de $f$ et dresser son tableau de variation.
  \item Démontrer que la droite $(\Delta)$ d'équation $y = 2 - x$ est asymptote à la courbe $(C)$ de la fonction $f$ et étudier la position de la droite $(\Delta)$ par rapport à cette courbe.
  \item Écrire l'équation de la tangente $(T)$ à la courbe $(C)$ de $f$ en $x = 0$.
  \item Étudier les branches infinies de la courbe $(C)$ de la fonction $f$.
  \item Construire dans le même repère, la courbe $(C)$ de $f$, la droite $(\Delta)$ et la tangente $(T)$.
  \item Calculer l'aire $\mathcal{A}(D)$ du domaine du plan limité par la courbe $(C)$, la droite $(\Delta)$ et les axes $x = \dfrac{3}{2}$ et $x = e$. On prendra $\ln 2 \approx 0{,}7$ ; $e \approx 2{,}7$.
\end{enumerate}

\end{document}
